\documentclass[border=12pt]{standalone}
\usepackage[utf8]{inputenc}
\usepackage[T1]{fontenc}
\usepackage[american, siunitx]{circuitikz}

\ctikzset{
    resistors/scale=0.85,
    diodes/scale=0.85,
    switches/scale=0.85
}

\begin{document}
\begin{circuitikz}[font=\sffamily\small]

    % =========================================================================
    % 1. BLOCO CENTRAL: ARDUINO UNO R3 (U1)
    % =========================================================================
    \draw[thick, fill=gray!5] (0, -4.2) rectangle (3.6, 4.2);
    \node[align=center, font=\sffamily\bfseries] at (1.8, 0.3) {U1\\Arduino\\Uno R3};
    \node[font=\sffamily\scriptsize, text=gray!70!black] at (1.8, -0.7) {(ATmega328P)};

    % Alimentacao do Arduino
    \draw (0, 3.4) -- (-0.8, 3.4) node[vcc]{+5V};
    \node[right, font=\ttfamily\footnotesize] at (0, 3.4) {5V};

    \draw (0, -3.4) -- (-0.8, -3.4) node[ground]{};
    \node[right, font=\ttfamily\footnotesize] at (0, -3.4) {GND};

    % Pinos de Entrada (Esquerda) com Net Labels
    \foreach \y/\pin/\lbl in {
        2.0/D9/BTN\_GREEN,
        0.8/D8/BTN\_RED,
        -0.4/D7/BTN\_BLUE,
        -1.6/D6/BTN\_YELLOW
    }{
        \node[right, font=\ttfamily\footnotesize] at (0, \y) {\pin};
        \draw (0, \y) -- (-1.2, \y) node[left, font=\ttfamily\scriptsize, fill=white, inner sep=1pt] {\lbl};
    }

    % =========================================================================
    % 2. BLOCO ESQUERDO: 4 BOTOES COM RESISTORES DE PULL-UP (1k)
    % =========================================================================
    % Barramento +5V superior com ponto de juncao no terminal de alimentacao
    \draw (-13.0, 3.4) -- (-4.6, 3.4);
    \node[vcc] at (-8.8, 3.4) {+5V};
    \node[circ] at (-8.8, 3.4) {};

    % Barramento GND inferior com ponto de juncao no terminal de terra
    \draw (-13.0, -2.4) -- (-4.6, -2.4);
    \node[ground] at (-8.8, -2.4) {};
    \node[circ] at (-8.8, -2.4) {};

    % Ramos dos 4 botoes (sem pontos nos cantos externos)
    \foreach \x/\rname/\swname/\lbl in {
        -13.0/R5/SW1/BTN\_GREEN,
        -10.2/R6/SW2/BTN\_RED,
        -7.4/R7/SW4/BTN\_BLUE,
        -4.6/R8/SW3/BTN\_YELLOW
    }{
        % Resistor de Pull-up (1 kOhm) - ponto apenas na derivacao central (-*)
        \draw (\x, 3.4) to[R, l_=\rname, a^=1<\kilo\ohm>, -*] (\x, 0.5);
        % Botao tatil (Pushbutton NA) - sem ponto na base
        \draw (\x, 0.5) to[push button, l_=\swname] (\x, -2.4);
        % Derivacao com Net Label
        \draw (\x, 0.5) -- ++(0.6, 0) node[right, font=\ttfamily\scriptsize] {\lbl};
    }

    % Pontos de juncao apenas nos ramos internos (R6/SW2 e R7/SW4)
    \foreach \x in {-10.2, -7.4}{
        \node[circ] at (\x, 3.4) {};
        \node[circ] at (\x, -2.4) {};
    }

    % =========================================================================
    % 3. BLOCO DIREITO SUPERIOR: 4 LEDs COM RESISTORES LIMITADORES (220R)
    % =========================================================================
    \foreach \y/\pin/\rname/\dname/\cor in {
        3.0/D13/R1/D1/Verde,
        1.8/D12/R2/D2/Vermelho,
        0.6/D11/R3/D3/Azul,
        -0.6/D10/R4/D4/Amarelo
    }{
        \node[left, font=\ttfamily\footnotesize] at (3.6, \y) {\pin};
        \draw (3.6, \y) -- (4.4, \y)
            to[R, l^=\rname, a_=220<\ohm>] (6.8, \y)
            to[leDo, l^={\dname\ (\cor)}] (9.4, \y)
            -- (10.2, \y);
    }

    % Barramento comum de GND dos LEDs
    \draw (10.2, 3.0) -- (10.2, -1.4) node[ground]{};
    \foreach \y in {1.8, 0.6, -0.6}{
        \node[circ] at (10.2, \y) {};
    }

    % =========================================================================
    % 4. BLOCO DIREITO INFERIOR: BUZZER PASSIVO (LS1)
    % =========================================================================
    \node[left, font=\ttfamily\footnotesize] at (3.6, -2.6) {D2};
    \draw (3.6, -2.6) -- (5.8, -2.6)
        to[buzzer, l^=LS1, a_={Buzzer Passivo}] (8.2, -2.6)
        -- (10.2, -2.6) -- (10.2, -3.4) node[ground]{};

\end{circuitikz}
\end{document}